\documentclass[11pt]{article}
\usepackage[margin=1in]{geometry}
\usepackage{amsmath,amssymb}
\usepackage{enumitem}
\usepackage{booktabs}
\usepackage{framed}

\newcommand{\indep}{{\bot\negthickspace\negthickspace\bot}}
\DeclareMathOperator{\E}{\mathbb{E}}

\pagestyle{plain}

\begin{document}

\begin{center}
{\Large \textbf{PS200C: Quantitative Methods}} \\[6pt]
{\large In-class Activity 7: Difference-in-Differences}
\end{center}

\vspace{4pt}
\noindent\textbf{Name:} \hrulefill

\vspace{12pt}

%%% QUESTION 1 %%%
\noindent\textbf{Question 1.}  Which of the following ``modern DID'' approaches can rescue your analysis if the \textbf{parallel trends} assumption is violated? (Circle all that apply.)

\medskip
\begin{enumerate}[label=(\alph*), itemsep=0.5em]
\item Sun \& Abraham (2021) cohort-specific event study
\item Callaway \& Sant'Anna (2021) group-time $ATT(g,t)$
\item Borusyak, Jaravel \& Spiess (2021) / Liu, Wang \& Xu (2022) imputation (\texttt{fect})
\item de Chaisemartin \& D'Haultf{\oe}uille (2020) convex-weighted clean 2$\times$2s
\item None of the above
\end{enumerate}

 \vspace{0.5em}
% \noindent What core problem do the modern DID approaches focus on?

\vspace{1cm}

%%% SETUP FOR Q2-Q4 %%%
\begin{framed}
\noindent\textbf{Setup for Q2--Q4.} In January 2018, Seattle imposed a sugary-drink tax.  Portland did not.  You have annual per-capita sales of sugary drinks (liters per person per year) in both cities, averaged over a short pre-window (2016--17) and a short post-window (2018--19):

\medskip
\begin{center}
\begin{tabular}{lcc}
\toprule
            & \textbf{Pre-tax (2016--17)} & \textbf{Post-tax (2018--19)} \\
\midrule
Seattle     &  100 liters  &  85 liters \\
Portland    &   95 liters  &  88 liters \\
\bottomrule
\end{tabular}
\end{center}
\end{framed}

\vspace{6pt}

%%% QUESTION 2 %%%
\noindent\textbf{Question 2: Compute the DID.}  Using the four means above, compute the difference-in-differences estimate of the effect of the tax on Seattle's per-capita sugary-drink sales.  Show your work.

\vspace{5cm}

\newpage

%%% QUESTION 3 %%%
\noindent\textbf{Question 3: Where might parallel trends fail here?}

\medskip
\begin{enumerate}[label=(\alph*), itemsep=0em, topsep=0pt]
\item Explain what parallel trends requires us to believe here.

\vspace{3.5cm}

\item Describe one concrete reason why parallel trends might fail.  \textit{Hint: think about something that could be changing during 2016--2019 that would influence both the sugary-drink-sales trend (absent the tax) and whether Seattle adopted the tax.}

\vspace{6cm}
\end{enumerate}

%%% QUESTION 4 %%%
\noindent\textbf{Question 4: Pre-trends check.}  Your collaborator plots five years of pre-tax data (2013--2017) and shows that Seattle and Portland's per-capita sales were declining at nearly the same rate over that window.  They conclude, ``Parallel trends is satisfied---we're good.''  Is that right? Why or why not?

\vspace{8cm}

\end{document}
